% ============================================================
% Paper 4 — A Dual Operator for Prime–Zero Coupling
%            and a Conditional Proof of Energy Asymmetry
%
% Author:  Ulrich Tehrani
% Version: v5 (April 2026)
% DOI:     10.5281/zenodo.19364703
% License: CC BY 4.0
%
% Changelog
% ---------
% v5 (April 2026)
%   Series consolidation (template v1.3/v1.4/v1.6):
%   - Notation and Conventions section added
%   - cmap + emergencystretch + frontmatter pagination
%   - Bibliography: Concept-DOIs, name-based keys,
%     Papers 4–6 added, Weil Suppl. precision
%   - Source hygiene: internal labels removed
%   - Footer: single-line format (v1.6)
%   - Assumption A(a)/A(b) cross-references clarified
%     (A(a) is settled by Hadamard throughout)
%
% v4 (April 2026)
%   Assumption A(a) correctly marked SETTLED (Hadamard);
%   Weyl-bridge Remark added; Theorem 5.3 retitled
%   CONDITIONAL/NUMERICAL; Call for Feedback added.
%
% v3 (March 2026)
%   Initial publication.
% ============================================================
\documentclass[11pt]{article}
\usepackage{cmap}            % ToUnicode CMaps for §, en/em-dashes etc.
\usepackage[T1]{fontenc}
\emergencystretch=1em        % avoid overfull hboxes from inline math
\usepackage[utf8]{inputenc}
\usepackage[a4paper,margin=1in]{geometry}
\usepackage{amsmath,amssymb,amsthm,mathtools}
\usepackage{xcolor}
\usepackage{hyperref}
\usepackage{enumitem}
\usepackage{setspace}
\usepackage{booktabs}
\usepackage{array}
\usepackage{graphicx}

\hypersetup{colorlinks=true,linkcolor=blue!50!black,urlcolor=blue!50!black,citecolor=blue!50!black}
\setlength{\parskip}{0.5em}
\setlength{\parindent}{0pt}
\onehalfspacing

% -- Theorem environments
\newtheorem{theorem}{Theorem}[section]
\newtheorem{proposition}[theorem]{Proposition}
\newtheorem{lemma}[theorem]{Lemma}
\newtheorem{corollary}[theorem]{Corollary}
\newtheorem{definition}[theorem]{Definition}
\newtheorem{assumption}[theorem]{Assumption}
\newtheorem{observation}[theorem]{Observation}
\newtheorem*{remark}{Remark}
\newtheorem{openproblem}[theorem]{Open Problem}

% Custom label-free environments
\newenvironment{model}[1]{\par\smallskip\noindent\textit{[Model: #1]}\quad}{\par\smallskip}
\newenvironment{heuristic}[1]{\par\smallskip\noindent\textit{[Heuristic: #1]}\quad}{\par\smallskip}

% -- Macros (Paper 3 base)
\newcommand{\Hstr}{H_{\mathrm{str}}}
\newcommand{\Hnull}{H_{\mathrm{null}}}
\newcommand{\Estr}{E_{\mathrm{str}}}
\newcommand{\Espec}{E_{\mathrm{spec}}}
\newcommand{\etaorig}{\eta_{\mathrm{orig}}}
\newcommand{\Tren}{T_{\mathrm{ren}}}
\newcommand{\Gun}{G^{\mathrm{un}}}
\newcommand{\Hlocal}{H_{\mathrm{local}}}
% -- Paper 4 specific macros
\newcommand{\Tt}{\widetilde{T}}
\newcommand{\Ceta}{C_{\eta}}
\newcommand{\CT}{C_{T}}
\newcommand{\kdom}{k_{\mathrm{dom}}}
\newcommand{\Wone}{W_1}
\newcommand{\Dburstvar}{\Delta_{\mathrm{Burst}}}
\newcommand{\Dcross}{\Delta_{\mathrm{Cross}}}
\newcommand{\Dstream}{\Delta_{\mathrm{Stream}}}

% -- Title
\title{\vspace{-1.5em}\bfseries A Dual Operator for Prime--Zero Coupling\\[0.4em]
\large and a Conditional Proof of Energy Asymmetry}
\author{Ulrich Tehrani}
\date{Research Note \textperiodcentered\ April 2026 \textperiodcentered\ v5}

\begin{document}
\maketitle

% -------------------------------------------------------
\begin{abstract}
We introduce the operator $\Tt := \Phi \circ \Phi^*$,
the dual of the loop operator $T = \Phi^* \circ \Phi$ constructed in Paper~3
of this series. Here $\Phi\colon\Hstr \to \Hnull$ maps a
finite-dimensional prime space into a finite-dimensional zero space via
Gaussian-damped prime resonance vectors.

\textbf{What is proved in this paper.}
We establish (\S\,2.3) the algebraic spectral identity
$\sigma(T)\setminus\{0\} = \sigma(\Tt)\setminus\{0\}$
by classical operator theory; the maximum discrepancy between shared
eigenvalues is below $10^{-15}$ (machine precision).
We prove (\S\,3.4) that $\Wone := \CT\cdot\Tt^+$ is self-adjoint.
$\Wone$ does not assume RH or the HP postulate; however, its normalization
constant $\CT$ is calibrated against the zero ordinates via OLS and should
not be presented as fully input-free.
$\Tt$ is explicit and deterministic; the $\gamma_k$ are external numerical
inputs (Definition~2.2), not conclusions of the model (\S\,3.4).
We prove (\S\,4.3, Lemma~4.3) the \textbf{Abel Summation Principle} (Lemma~M3):
for monotone decreasing weights $f_i$ and a sequence $g_i$ with bounded
partial sums, $|\sum_i f_i g_i| \leq M f_1$.
We prove (\S\,4.4, Theorem~4.4) the \textbf{Prime Exponential Sum Bound}
(Theorem~EXPLICIT): for fixed $\gamma > 0$,
\[
M_k(\kappa) := \max_{x \leq \kappa}\Bigl|\sum_{p \leq x}\sin(\gamma_k\log p)\Bigr|
= O\!\left(\frac{\pi(\kappa)}{\gamma_k}\right)
\]
using only the Prime Number Theorem (no Riemann Hypothesis).
We prove (\S\,2.4, Proposition~2.11) that $\Tt$ is \textbf{not} a
Hilbert--P\'olya operator: the correlation
$r_2 = \mathrm{corr}(\omega_j,\gamma_{k(j)})$
falls from $0.50$ to $0.16$ for $\kappa\in\{23,\ldots,503\}$.

\textbf{What is numerical.}
The eigenvectors of $\Tt$ localize at zero ordinates $\gamma_k$
with degrees $0.45$--$0.99$ [Numerical, \S\,3.2].
The correlation $r_1 = \mathrm{corr}(\mu_j,1/\gamma_{k(j)}) = 0.950$
($\kappa=53$) [Numerical, \S\,3.1].
The energy asymmetry $\etaorig > 0$ holds for all tested $\kappa\leq 1009$
[Numerical, \S\,5.1].

\textbf{What is conditional.}
Under Assumption~A --- a condition on the line
$\mathrm{Re}(s)=1$, distinct from the Riemann Hypothesis:
$(a)$~$\zeta(1+in\gamma_k)\neq 0$ for all $n\geq 1$, $k\leq N$
(this follows directly from Hadamard--de la Vall\'ee Poussin (1896),
which establishes $\zeta(1+it)\neq 0$ for all real $t\neq 0$;
hence not genuinely open), and
$(b)$~joint equidistribution of prime log-phases (\S\,6, Open Problem~6.1)
(the genuinely open condition)
--- we prove (\S\,5.3, Theorem~5.3(i)(ii)) that the normalised cross/stream
averages vanish [CONDITIONAL on Assumption~A].
We establish numerically (Theorem~5.3(iii)(iv)) that
$\Delta(\kappa)\geq\Dburstvar\approx 3.11 > 0$ for all tested $\kappa > 7$
[NUMERICAL lower bound, $\kappa$-invariant] and
$\etaorig(\kappa)\to\eta_\infty\approx 0.81 > 0$ [NUMERICAL].

\textbf{What remains open.}
The analytic proof of $\etaorig > 0$ without the above assumption
(Open Problem~6.2); whether a circularity-free spectral function
$f(\Tt)$ exists with spectrum converging to $\{\gamma_k\}$
(Open Problem~6.3); and the arithmetic origin of
$\Ceta \approx 0.39$ (Open Problem~6.4).

\textbf{Structure.}
The formal core in \S\S\,2--5 is non-circular and proof-complete.
Sections~6 and~7 contain open problems and a geometric interpretation
(clearly marked, not used in proofs).
No use of the Riemann Hypothesis, the GUE conjecture, the Montgomery
pair correlation conjecture, or the Hilbert--P\'olya postulate is made
anywhere in \S\S\,2--5.
\end{abstract}

\vspace{2em}

% --------------------------------------------------------
% Notation and Conventions  (template invariant, v1.3)
% --------------------------------------------------------
\section*{Notation and Conventions}
\label{sec:notation}

\emph{Critical line.}
Throughout this series, the \emph{critical line} means
$\Re(s)=\tfrac12$, parameterised by $s=\tfrac12+i\gamma$
with $\gamma\in\mathbb{R}$.

\medskip
\emph{Global parameters.}
All numerical results use the normative parameter set:
prime cutoff $\kappa=53$ ($\pi(\kappa)=16$ active primes),
Gaussian regularisation parameter $\varepsilon=0.05$,
$N=100$ Riemann zero ordinates, and evaluation point
$\sigma_0=\tfrac12$.

\medskip
\emph{Main objects.}
The encoding map $\Phi\colon\Hstr\to\Hnull$
assigns to each prime $p\leq\kappa$ a resonance vector
$a_p=(e^{-\varepsilon^2\gamma_k^2/2}\sin(\gamma_k\log p))_{k=1}^N$.
The loop operator $T=\Phi^*\!\circ\Phi$ acts on the prime
space $\Hstr$; its dual, the Tehrani operator
$\Tt=\Phi\circ\Phi^*$, acts on the zero space $\Hnull$
and is the central object of this paper.
The energy asymmetry $\etaorig=1-c^T\Gun c/\Estr$ measures
the gap between structural and spectral energy.
Two arithmetic constants govern spectral scaling:
$\Ceta\approx 0.39$ (renormalised eigenvalue growth) and
$\CT$ (OLS slope of $\mu_j$ vs.\ $1/\gamma_{k(j)}$).

\medskip
\emph{Not to be confused with.}
$\Ceta\neq\CT$: $\Ceta$ governs
$\lambda_{\max,\mathrm{ren}}\sim\Ceta\cdot\pi(\kappa)$,
while $\CT$ governs $\mu_j\sim\CT/\gamma_{k(j)}$.
The inverse operator $W_1=\CT\cdot\Tt^+$ is \emph{not}
a Hilbert--P\'olya operator; its HP-correlation $r_2$
decays to $0.16$ with growing $\kappa$.

\newpage

% -------------------------------------------------------
\section{Introduction}

\subsection{Background and Motivation}

This paper is the fourth in a series studying the curvature of the Riemann
zeta function from a geometric and operator-theoretic perspective.
A reader familiar with Paper~3 \cite{Paper3} may begin at \S\,2;
the present introduction provides series context and an explicit
reader contract for all results.
The formal content in \S\S\,2--5 is non-circular and proof-complete.

\textbf{Paper~1} \cite{Paper1} introduced the local curvature function
\[
H_\xi(\sigma,\kappa) := \sum_{p\leq\kappa} V_p(\sigma),
\]
where $V_p(\sigma) = \tfrac{d^2}{d\sigma^2}\log\|M_\sigma f_{\kappa,p}\|^2 > 0$.
A key finding was the divergence
$H_{\mathrm{local}}(\tfrac{1}{2},\kappa)\sim 8(\log\kappa)^2\to\infty$
at $\sigma=\tfrac{1}{2}$, contrasted with
$H_{\mathrm{local}}(\sigma,\kappa)\to C(\sigma)<\infty$
for all $\sigma>\tfrac{1}{2}$.

\textbf{Paper~2} \cite{Paper2} constructed the Weil functional and established
$W(g^*\ast\check{g}^*) = Z(g^*) - H_{\mathrm{local}}(\sigma,\kappa) + O(\varepsilon)$,
connecting local curvature to the Weil positivity criterion.

\textbf{Paper~3} \cite{Paper3} made the geometric structure explicit.
Given a prime cutoff $\kappa$ and $N$ zero ordinates $\gamma_k$,
two finite-dimensional real Hilbert spaces $\Hstr$ (over primes) and
$\Hnull$ (over zeros) are connected by a linear map $\Phi$.
The loop operator $T = \Phi^*\circ\Phi$ is self-adjoint and positive
semi-definite.
The energy asymmetry
$\etaorig = \Delta/\Estr \in [0.598, 0.704]$
(benchmark $\kappa=53$, $\varepsilon=0.05$, $N=100$) was established
numerically; the structural energy $\Estr\to 6.1$ as $\kappa\to\infty$
(analytically explained by absolute convergence of $\sum_p(\log p)^2/p^2$).

\textbf{The present paper} introduces and analyzes
$\Tt = \Phi\circ\Phi^*$, the dual of $T$, acting on the zero space
$\Hnull$.
The operator emerges canonically from the same $\Phi$ already constructed
in Paper~3, and requires no new postulate.

The main conditional result of Paper~4
is a \textbf{conditional analytic proof} that $\etaorig > 0$: under Assumption~A --- two conditions
on $\mathrm{Re}(s)=1$ --- the energy asymmetry persists in the limit
$\kappa\to\infty$.
Both conditions concern $\mathrm{Re}(s)=1$, distinct from RH which concerns
$\mathrm{Re}(s)=\tfrac{1}{2}$.
Non-vanishing on $\mathrm{Re}(s)=1$ is established for all real $t\neq 0$
by Hadamard and de la Vall\'ee Poussin (1896); this covers the values
$t = n\gamma_k$ directly.
The genuinely open condition is the equidistribution of
$\{\gamma_k\log p\bmod 2\pi\}$, which is the main open problem of
this paper.

\subsection{Non-Circularity Statement}

\emph{The following holds unconditionally throughout \S\S\,2--5:}

\textbf{No result in this paper assumes:}
\begin{enumerate}[label=(\roman*)]
\item The Riemann Hypothesis;
\item The GUE conjecture;
\item The Montgomery pair correlation conjecture;
\item The Hilbert--P\'olya postulate.
\end{enumerate}

The operator $\Tt = \Phi\circ\Phi^*$ is constructed explicitly from
the prime resonance vectors $\{a_p\}_{p\leq\kappa}$ (defined in \S\,2.1).
The zero ordinates $\gamma_k$ appear as \emph{inputs} to the construction
of $\Phi$, not as conclusions.
In Theorem~5.3 (\S\,5.3), Assumption~A is explicitly
labelled CONDITIONAL; it is a condition on $\mathrm{Re}(s)=1$, distinct
from RH which concerns $\mathrm{Re}(s)=\tfrac{1}{2}$.
Its logical relationship to RH is not analysed in this paper.

In particular, $\Tt = \Phi\circ\Phi^*$ is a deterministic,
explicitly computable finite-dimensional operator --- not a random matrix
or a statistical ensemble.
No GUE conjecture, no Montgomery pair correlation conjecture,
and no random-matrix-theoretic assumption is used or implied anywhere
in \S\S\,2--5.
The structural similarity between eigenvector localization in $\Tt$ and
GUE level statistics is a numerical observation, not a postulate.

\emph{Note on $\Wone$:}
The inverse operator $\Wone = \CT\cdot\Tt^+$ does not assume RH or the
HP postulate; however, its normalization constant $\CT$ is calibrated
against the zero ordinates via OLS (\S\,3.3).
$\Wone$ is therefore not fully input-free.
$\Tt$ is deterministic and explicitly computable; the $\gamma_k$ are
external numerical inputs, not conclusions of the model.

\subsection{Imported Results from Papers 1--3}

The following results from earlier papers are used in this paper,
exclusively in the indicated roles.

\begin{center}
\small
\renewcommand{\arraystretch}{1.3}
\begin{tabular}{lp{4.5cm}p{2.8cm}p{3cm}}
\toprule
Label & Result & Status & Used where \\
\midrule
(R1) & $H_{\mathrm{local}}(\tfrac{1}{2},\kappa)\sim 8(\log\kappa)^2$
  & Proved [Paper~1, \S\,2] & \S\,1.1, \S\,7 (motivation only) \\
(R2) & Weil bridge
  & Proved [Paper~2, \S\,3] & \S\,1.1, \S\,7 (motivation only) \\
(R3) & $T=\Phi^*\Phi$, $\etaorig\in[0.598,0.704]$, normative $c_p$
  & Proved/Numerical [Paper~3] & \S\,2.1, \S\,5 \\
(R4) & $\Estr(\infty)\approx 6.1$ (convergent)
  & Numerical [Paper~3, \S\,5] & \S\,5.1 (context) \\
\bottomrule
\end{tabular}
\end{center}

Results (R1) and (R2) are used for motivation only.
Results (R3) and (R4) are recorded for completeness;
the formal core in \S\S\,2--5 depends only on the definitions in \S\,2.1,
not on $\eta$-values from Paper~3.

\subsection{What This Paper Proves and What It Imports}

\emph{Reader Contract.}

\begin{center}
\renewcommand{\arraystretch}{1.3}
\begin{tabular}{lll}
\toprule
Content & Status & Where \\
\midrule
$\Tt = \Phi\circ\Phi^*$ definition & Proved here & \S\,2.2 \\
Self-adjointness and PSD of $\Tt$ & Proved here & \S\,2.2 \\
Algebraic spectral identity $\lambda_j(T)=\mu_j(\Tt)$ & Proved here & \S\,2.3 \\
$\Tt$ is NOT a Hilbert--P\'olya operator & Num.\ negative result & \S\,2.4 \\
Lemma M3 (Abel Summation Principle) & Proved here & \S\,4.3 \\
Theorem EXPLICIT: $M_k(\kappa)=O(\pi(\kappa)/\gamma_k)$ & Proved here & \S\,4.4 \\
$\Wone = \CT\cdot\Tt^+$ self-adjoint & Proved here & \S\,3.4 \\
$\Wone$: no RH/HP; $\CT$ OLS-calibrated & Proved here & \S\,3.4 \\
$\mu_j \sim \CT/\gamma_{k(j)}$, $r_1=0.950$ & \textbf{Numerical} & \S\,3.1 \\
Eigenvector localization $0.45$--$0.99$ & \textbf{Numerical} & \S\,3.2 \\
$\etaorig > 0$ for $\kappa\leq 1009$ & \textbf{Numerical} & \S\,5.1 \\
$\Dburstvar \approx 3.11 > 0$, $\kappa$-invariant & \textbf{Numerical} & \S\,5.2 \\
$\Wone$ is NOT an asymptotic HP-operator & \textbf{Numerical} & \S\,3.4 \\
Theorem~5.3(i)(ii) (equidist.\ vanishing) & \textbf{Conditional} & \S\,5.3 \\
Theorem~5.3(iii): $\Delta(\kappa)\geq\Dburstvar > 0$
  & \textbf{Numerical} & \S\,5.3 \\
Theorem~5.3(iv): $\etaorig(\kappa)\to\eta_\infty\approx 0.81 > 0$
  & \textbf{Numerical} & \S\,5.3 \\
$\etaorig > 0$ analytically & \textbf{Open} & \S\,6 \\
Assumption~A(a): $\zeta(1+in\gamma_k)\neq 0$
  & \textbf{Settled} (Hadamard 1896) & \S\,6.1 \\
Assumption~A(b): equidistribution $\{\gamma_k\log p\bmod 2\pi\}$
  & \textbf{Open} & \S\,6.1 \\
\bottomrule
\end{tabular}
\end{center}

\subsection{Relation to Existing Work}

\textbf{Hilbert--P\'olya setting.}
A Hilbert--P\'olya operator $\mathcal{H}$ would satisfy
$\mathcal{H}\psi_k = \gamma_k\psi_k$ ---
the zero ordinates appear directly as eigenvalues.
The Tehrani operator $\Tt$ has eigenvalues $\mu_j\sim\CT/\gamma_{k(j)}$
(numerically); the zero ordinates appear as \emph{localization centers}
of eigenvectors, not as eigenvalues.
This is a weaker but constructive relationship.
The inverse $\Wone = \CT\cdot\Tt^+$
is a natural candidate for an HP-operator,
but it fails (\S\,2.4): the correlation
$r_2 = \mathrm{corr}(\omega_j,\gamma_{k(j)})$ falls to $0.16$.
We use the standard formulation of the Hilbert--P\'olya setting
--- see e.g.\ the survey by Bombieri \cite{Bombieri} ---
without adopting it as a postulate.

\begin{quote}
\textit{Terminology (used throughout this paper):}
We call $\Tt$ a ``resonance operator'' to express the fact that
its eigenvectors concentrate around prescribed arithmetic centers
$\gamma_k$.
This is \emph{paper-specific terminology}: in spectral theory,
``resonance'' usually refers to scattering poles or quasimodes.
The term is not used in the scattering-theoretic sense here.
\end{quote}

\textbf{Positioning relative to four established programs.}
Any operator-theoretic study of zeta zeros is naturally compared
with the following programs.

(a)~\emph{Berry--Keating \cite{BerryKeating}.}
The $xp$-quantisation program seeks a self-adjoint operator
whose eigenvalues are the zero ordinates $\gamma_k$.
Paper~4 operates in the opposite direction: $\gamma_k$ appear
as eigenvector localization centers, not eigenvalues, and $\Tt$
is a finite-dimensional Gram-type operator rather than a
semiclassical Hamiltonian.

(b)~\emph{Connes \cite{Connes}.}
Connes gives an adelic trace-formula interpretation in which
critical zeros appear as an absorption spectrum.
The present construction is finite-dimensional, built directly
from prime resonance vectors, and involves no adelic or
noncommutative-geometric machinery.
No formal connection between $\Tt$ and the Connes trace formula
is established in this paper.

(c)~\emph{Sierra--Townsend \cite{SierraTownsend}.}
Their Landau-level model provides a truncated Hamiltonian
realization connected to Connes' absorption-spectrum picture.
Paper~4's $\Tt$ is not a Hamiltonian but a dual encoding operator;
the spectral mechanism is different.

(d)~\emph{Bost--Connes \cite{BostConnes}.}
The Bost--Connes system is a $C^*$-dynamical system with zeta
partition function, operating in a quantum-statistical framework.
Paper~4's construction is purely finite-dimensional and
linear-algebraic, with no KMS-state or phase-transition structure.

\textbf{Goldston--Gonek \cite{GoldstonGonek}.}
Their exponential sum formula provides an independent confirmation of
$M_k(\kappa) = O(\pi(\kappa)/\gamma_k)$ proved in \S\,4.4
(see Remark~4.7).

\textbf{Tur\'an \cite{Turan}.}
The closest historical precedent for sums of the form
$\sum_p e^{-i\tau\log p}$ is Tur\'an's
program on approximative Dirichlet polynomials.
The estimate proved in Theorem~\ref{thm:EXPLICIT}
uses standard PNT/Abel-summation methodology;
the novelty lies not in the bare oscillatory-sum bound itself,
but in how it enters the conditional energy-asymmetry
argument of \S\,5.

\textbf{Paper~3 \cite{Paper3}.}
The present paper is the direct sequel; \S\,2.1 recollects all necessary
objects.

\subsection{Outline}

Section~2 defines $\Tt$, proves the algebraic spectral identity,
and establishes that $\Tt$ is not a Hilbert--P\'olya operator.
Section~3 presents the spectral analysis of $\Tt$: eigenvalue structure,
eigenvector localization, and the two arithmetic constants
$\Ceta$ and $\CT$.
Section~4 develops the monotonicity mechanism: Abel summation (Lemma~M3)
and the prime exponential sum bound (Theorem~EXPLICIT).
Section~5 collects results on $\etaorig > 0$: numerical confirmation,
the $\Delta$-decomposition, and the conditional theorem.
Section~6 states the open problems.
Section~7 presents the Fisher-information geometric interpretation
(all statements in \S\,7 are explicitly marked as geometric model
or heuristic).

% -------------------------------------------------------
\section{The Tehrani Operator}

\subsection{Recollection: $\Hstr$, $\Hnull$, $\Phi$, $T$}

We fix throughout:
\begin{itemize}
\item A prime cutoff $\kappa > 1$, with active prime set
  $S(\kappa) := \{p\text{ prime}: p\leq\kappa\}$
  and $\pi(\kappa) := |S(\kappa)|$.
\item The first $N$ positive imaginary parts of nontrivial zeros of
  $\zeta(s)$, ordered $0 < \gamma_1 \leq \gamma_2 \leq \cdots \leq \gamma_N$.
\item A Gaussian damping parameter $\varepsilon > 0$.
\end{itemize}

\textbf{Normative parameters:} $\kappa=53$, $\varepsilon=0.05$,
$N=100$, $\sigma=\tfrac{1}{2}$.

\begin{definition}[$\Hstr$ and $\Hnull$, from Paper~3]
\[
\Hstr := \ell^2(S(\kappa)),\quad \dim\Hstr = \pi(\kappa).
\qquad
\Hnull := \ell^2(\{\gamma_k\}_{k=1}^N),\quad \dim\Hnull = N.
\]
Both spaces carry standard inner products.
We work throughout with real vectors and real operators.
\end{definition}

\begin{definition}[Resonance vectors, from Paper~3]
For each prime $p\in S(\kappa)$, the \emph{resonance vector}
$a_p\in\Hnull$ is defined by
\[
a_p := \bigl(e^{-\varepsilon^2\gamma_k^2/2}\sin(\gamma_k\log p)\bigr)_{k=1}^N.
\]
The Gaussian envelope ensures $a_p\in\ell^2(\{\gamma_k\})$ uniformly
in $\kappa$.
\end{definition}

\begin{definition}[Linear map $\Phi$, from Paper~3]
\[
\Phi\colon\Hstr\to\Hnull,\quad
\Phi(e_p) = a_p,\quad
\Phi(c) = \sum_{p\leq\kappa}c_p a_p.
\]
\end{definition}

\begin{definition}[Loop operator $T$, from Paper~3]
\[
T := \Phi^*\circ\Phi\colon\Hstr\to\Hstr,\quad
T_{pq} = \langle a_p,a_q\rangle_{\Hnull} = G^{\mathrm{un}}_{pq}.
\]
$T$ is self-adjoint, positive semi-definite, and identical to the Gram
matrix $\Gun$ of the resonance vectors.
\end{definition}

\begin{definition}[Canonical weight vector, from Paper~3]
The normative weight vector $c\in\Hstr$ has components
$c_p := \sqrt{V_p(\tfrac{1}{2})}$,
where $V_p(\sigma) = \tfrac{d^2}{d\sigma^2}\log\|M_\sigma f_{\kappa,p}\|^2$.
The exact formula is
$c_p = \sqrt{4(\log p)^2(2p-1)/(p(p-1)^2)}$.
The leading-term ($K=1$) approximation $c_p \approx 4\log(p)/\sqrt{p}$
is asymptotically close but numerically differs for small primes;
the exact formula is used throughout this paper and in all verification scripts.

Reference values ($\varepsilon=0.05$, $\sigma=\tfrac{1}{2}$):
$c_2\approx 1.698$, $c_3\approx 1.418$, $c_5\approx 1.080$,
$c_7\approx 0.884$.
\end{definition}

\begin{definition}[Energy quantities, from Paper~3]
\[
\Estr := \sum_{p\leq\kappa}c_p^2\|a_p\|^2 = c^T Dc,
\qquad
\Espec := \Bigl\|\sum_{p\leq\kappa}c_p a_p\Bigr\|^2 = \|\Phi c\|^2,
\]
where $D := \mathrm{diag}(\|a_p\|^2)_{p\in S(\kappa)}$.
The energy asymmetry is
\[
\Delta := \Estr - \Espec \geq 0\quad\text{(numerically)},
\qquad
\etaorig := \frac{\Delta}{\Estr}
= 1 - \frac{c^T\Gun c}{\Estr}.
\]
\end{definition}

\subsection{Definition of the Tehrani Operator $\widetilde{T}$}

\begin{definition}[Tehrani Operator]
\begin{equation}
\boxed{\Tt := \Phi\circ\Phi^*\colon\Hnull\to\Hnull.}
\end{equation}

\textbf{Matrix representation.}
In the standard basis $\{f_k\}_{k=1}^N$ of $\Hnull$:
\[
\Tt_{kl}
= \sum_{p\leq\kappa}(a_p)_k(a_p)_l
= \sum_{p\leq\kappa}
e^{-\varepsilon^2\gamma_k^2/2}\sin(\gamma_k\log p)
\cdot
e^{-\varepsilon^2\gamma_l^2/2}\sin(\gamma_l\log p).
\]
\end{definition}

\begin{proposition}[Properties of $\Tt$]\label{prop:Tt-properties}
The operator $\Tt$ satisfies:
\begin{enumerate}[label=(\roman*)]
	\item \textbf{Self-adjointness:}
	$\Tt^* = (\Phi\Phi^*)^* = \Phi\Phi^* = \Tt$. $\checkmark$
	\item \textbf{Positive semi-definiteness:}
	For all $u\in\Hnull$,
	$\langle\Tt u,u\rangle = \|\Phi^* u\|^2 \geq 0$. $\checkmark$
	\item \textbf{Rank:}
	$\mathrm{rank}(\Tt) = \mathrm{rank}(\Phi) = \pi(\kappa)$. $\checkmark$
	\item \textbf{Explicit construction:}
	$\Tt$ is built directly from prime resonance vectors $\{a_p\}$.
	The zero ordinates $\gamma_k$ are numerical inputs to the
	resonance vectors (Definition~2.2) and hence to $\Tt$.
	Non-circularity means: no RH, GUE, or HP assumption is used;
	the $\gamma_k$ are external inputs, not conclusions of the model.
\end{enumerate}
\emph{Proof.}
Parts (i) and (ii) are immediate from $\Tt=\Phi\Phi^*$ using standard
properties of adjoints.
Part (iii) follows from
$\mathrm{rank}(AB)\leq\min(\mathrm{rank}(A),\mathrm{rank}(B))$
and $\pi(\kappa)\leq N$ in all normative parameter choices. $\square$
\end{proposition}

\subsection{Algebraic Spectral Identity}

\begin{theorem}[Spectral Identity — PROVED]\label{thm:spectral-identity}
\[
\sigma(T)\setminus\{0\} = \sigma(\Tt)\setminus\{0\}.
\]
The shared nonzero eigenvalues satisfy:
\[
\max_{j=1}^{\pi(\kappa)}|\lambda_j(T) - \mu_j(\Tt)| < 10^{-15}
\quad\text{(machine precision, verified numerically)}.
\]
\emph{Proof.}
Classical result (see e.g.\ \cite{IwaniecKowalski}, Appendix~A,
or any standard linear algebra text): for any $A,B$,
$\sigma(AB)\setminus\{0\} = \sigma(BA)\setminus\{0\}$.
Applied with $A=\Phi^*$, $B=\Phi$: $T=\Phi^*\Phi$ and
$\Tt=\Phi\Phi^*$, so $\sigma(T)\setminus\{0\}=\sigma(\Tt)\setminus\{0\}$.
(For finite-dimensional matrices, Sylvester's determinant identity:
$\det(\lambda I - AB) = \lambda^{n-m}\det(\lambda I-BA)$
for $A\in\mathbb{R}^{m\times n}$, $B\in\mathbb{R}^{n\times m}$.) $\square$
\end{theorem}

\begin{remark}
The identity $\lambda_j(T)=\mu_j(\Tt)$ is an \emph{algebraic} theorem
--- it holds in exact arithmetic for any choice of $\kappa$, $\varepsilon$,
and $N$.
The eigenvalues $\lambda_j = \mu_j$ are not the zero ordinates $\gamma_k$;
they are constructive spectral quantities arising from $\Phi$.
\end{remark}

\subsection{$\Tt$ is NOT a Hilbert--P\'olya Operator}

This section records a negative result.

\begin{proposition}[Numerical Negative Result]\label{prop:not-HP}
A Hilbert--P\'olya operator $\mathcal{H}$ would satisfy
$\mathcal{H}\psi_k = \gamma_k\psi_k$ for all $k$,
i.e., the zero ordinates $\gamma_k$ appear directly as eigenvalues.
$\Tt$ does \textbf{not} have this property:
\begin{enumerate}[label=(\alph*)]
\item \emph{[Numerical]} The eigenvalues of $\Tt$ satisfy
  $\mu_j\sim\CT/\gamma_{k(j)}$ --- they are \emph{inversely} proportional
  to zero ordinates (\S\,3.1).
\item \emph{[Numerical]} The zero ordinates $\gamma_k$ appear as
  \emph{localization centers} of eigenvectors $v_j$ (\S\,3.2),
  not as eigenvalues.
\item \emph{[Numerical]} The inverse operator $\Wone := \CT\cdot\Tt^+$
  has eigenvalues $\omega_j = \CT/\mu_j \approx \gamma_{k(j)}$
  for small $j$, but the correlation
  $r_2 := \mathrm{corr}(\omega_j, \gamma_{k(j)})$
  falls from $0.502$ (at $\kappa=23$) to $0.160$ (at $\kappa=503$).
  $\Wone$ is therefore \textbf{not} an asymptotic HP-operator.
\end{enumerate}
\emph{The fundamental distinction:}
\[
\text{HP-operator: eigenvalues} = \gamma_k.
\qquad
\Tt\text{: eigenvector localization centers} = \gamma_k.
\]
$\Tt$ is a \textbf{resonance operator}: it encodes the $\gamma_k$
as mode centers, not as spectral values.

\emph{Proof of (a)(b)(c).}
Numerical; see \S\,3.1, \S\,3.2, \S\,3.4 for detailed data. $\square$
\end{proposition}

\begin{remark}[Structural obstruction]
The Gaussian damping $e^{-\varepsilon^2\gamma_k^2/2}$ in $a_p$ causes
$\mu_j$ to decay super-exponentially (exponent $\alpha\approx -4.88$),
far faster than the $1/\gamma$-law (exponent $-0.56$).
Inverting $\Tt$ to form $\Wone$ catastrophically amplifies the smallest
eigenvalues.
For $j=16$ at $\kappa=53$:
$\omega_{16}\approx 7\times 10^6$ against $\gamma_{16}\approx 67$.
This is the structural reason $\Wone$ fails as an HP-operator
(see \S\,3.4 and Open Problem~6.3).
\end{remark}

% -------------------------------------------------------
\section{Spectral Analysis of $\widetilde{T}$}

\subsection{Eigenvalue Structure: $\mu_j \sim \CT/\gamma_{k(j)}$}

\textbf{Numerical Result NR~3.1} (Eigenvalue--Zero Ordinate Correlation).
\emph{[Status: NUMERICAL --- $\kappa=53$, $\varepsilon=0.05$, $N=100$]}

\emph{All numerical results in this paper are reproducible using
the verification scripts at \url{https://github.com/utehrani/analysislab-nt}
(public repository, March 2026) \cite{Code}.
All computations use $N=100$ zero ordinates computed via \texttt{mpmath}
at default precision (approximately 15 significant digits) and normative
parameters $\kappa=53$, $\varepsilon=0.05$, $\sigma=\tfrac{1}{2}$
unless otherwise stated.}

Let $\kdom(j) := \arg\max_k|(v_j)_k|$ be the dominant zero index
for eigenvector $v_j$. Then:

\begin{center}
\begin{tabular}{ccccc}
\toprule
$j$ & $\mu_j(\Tt)$ & $\kdom$ & $\gamma_{\kdom}$ & Localization \\
\midrule
1 & 4.219803 & 1 & 14.1347 & 0.985 \\
2 & 2.565694 & 3 & 25.0109 & 0.741 \\
3 & 1.744649 & 2 & 21.0220 & 0.745 \\
4 & 0.986492 & 4 & 30.4249 & 0.854 \\
5 & 0.314199 & 5 & 32.9351 & 0.794 \\
6 & 0.217900 & 6 & 37.5862 & 0.885 \\
\bottomrule
\end{tabular}
\end{center}

The OLS correlation between $\mu_j$ and $1/\gamma_{k(j)}$ is
\[
r_1 = \mathrm{corr}\bigl(\mu_j,\,1/\gamma_{k(j)}\bigr) = 0.950
\quad(\kappa=53).
\]
For $\kappa=503$: $r_1 = 0.859$ (weakening as $\kappa$ grows,
due to Gaussian damping).

\begin{remark}
The dominant index satisfies $\kdom(1) = 1$ for all tested
$\kappa\in\{23,53,101,199,503,1009\}$:
the first eigenvector of $\Tt$ is always associated with
$\gamma_1 = 14.1347$.
\end{remark}

\begin{remark}[Cauchy--Schwarz warning]
Note that $r_1 = 0.950$ coexists with $r_2 = 0.502$ (\S\,3.4).
These measure different things: $r_1$ tests monotone ordering;
$r_2$ tests value-level approximation.
The distinction between ordering and value accuracy is crucial for
the HP assessment in \S\,2.4.
\end{remark}

\subsection{Eigenvector Localization}

\textbf{Numerical Result NR~3.2} (Eigenvector Localization).
\emph{[Status: NUMERICAL --- $\kappa=53$, $\varepsilon=0.05$, $N=100$]}

For each eigenvector $v_j$ of $\Tt$ (normalized: $\|v_j\|=1$),
the localization is measured by $L_j := |(v_j)_{\kdom(j)}|$.
Observed range: $L_j\in[0.45,0.99]$ for $j=1,\ldots,14$.

For comparison, a uniformly random unit vector in $\mathbb{R}^{100}$
has typical entry magnitude $1/\sqrt{100} = 0.10$.
The observed localization is $4.5$--$9.9$ times the random baseline.

\emph{[Geometric interpretation, not used in any proof]:}
The localization structure suggests that $\Tt$ `encodes'
the ordering of zero ordinates in a distributional sense:
each mode $v_j$ is concentrated near the zero $\gamma_{k(j)}$
in the zero-space $\Hnull$.
This is the precise content of calling $\Tt$ a resonance operator.

\subsection{The Two Arithmetic Constants $\Ceta$ and $\CT$}

\begin{definition}[Arithmetic Constants]\label{def:constants}
\[
\Ceta := \lim_{\kappa\to\infty}\frac{\lambda_{\max}(\Tren)}{\pi(\kappa)},
\qquad\Ceta\approx 0.39,
\]
where $\Tren := D^{-1/2}TD^{-1/2}$.
\[
\CT(\kappa) := \text{OLS-slope of } \mu_j\sim\CT/\gamma_{k(j)},
\qquad \CT/\pi(\kappa) \in \{2.49,\,2.21,\,2.17,\,1.94,\,1.89\}
\]
for $\kappa\in\{23,53,101,199,503\}$ (weakly decreasing).

\[
\boxed{\Ceta \neq \CT}
\]
\emph{These are two distinct arithmetic quantities. Always use subscripts.}
\end{definition}

The constant $\Ceta\approx 0.39$ is $\kappa$-independent;
$\CT$ grows with $\kappa$ (numerically observed to grow approximately as $2\pi(\kappa)$).

\textbf{Numerical Result NR~3.3} ($\lambda_{\max,\mathrm{ren}}$ growth).
\emph{[Status: NUMERICAL]}

\begin{center}
\begin{tabular}{cccc}
\toprule
$\kappa$ & $\pi(\kappa)$ & $\lambda_{\max,\mathrm{ren}}$ & Ratio \\
\midrule
23  & 9   & 3.59  & 0.399 \\
53  & 16  & 6.38  & 0.399 \\
101 & 26  & 10.35 & 0.398 \\
199 & 46  & 17.76 & 0.386 \\
503 & 96  & 37.38 & 0.389 \\
1009 & 169 & 65.16 & 0.386 \\
\bottomrule
\end{tabular}
\end{center}

Variance of the ratio: $\sigma^2 < 0.0003$ --- extremely stable.
The universal bound $\lambda_{\max,\mathrm{ren}} < 1$
(an earlier conjecture in this series, now falsified) is numerically falsified:
$\lambda_{\max,\mathrm{ren}}\to\infty$.

\begin{remark}[Analytical status of $\Ceta$]
The value $\Ceta\approx 0.39$ is not analytically explained.
A candidate formula $\Ceta = 1/I_N(\varepsilon)$ is \textbf{falsified}:
it gives $1/I_{100}(0.05)\approx 0.73\neq 0.39$.
The correct representation involves a harmonic mean
(analytically open; see Open Problem~6.4 and \S\,6).
\end{remark}

\subsection{The Inverse Operator $\Wone = \CT\cdot\Tt^+$}

\begin{definition}[Inverse Operator]
$\Wone := \CT\cdot\Tt^+$,
where $\Tt^+$ denotes the Moore--Penrose pseudoinverse of $\Tt$
and $\CT$ is the OLS constant from Definition~\ref{def:constants}.
\end{definition}

\begin{proposition}[Self-Adjointness of $\Wone$ --- PROVED]\label{prop:W1-sa}
$\Wone$ is self-adjoint.

\emph{Proof.}
$\Tt$ is self-adjoint (Proposition~\ref{prop:Tt-properties}(i)).
For any self-adjoint operator, the Moore--Penrose pseudoinverse is also
self-adjoint (standard result in linear algebra).
$\CT$ is a real scalar.
Therefore $\Wone = \CT\cdot\Tt^+$ is self-adjoint. $\square$
\end{proposition}

Numerically: $\|\Wone - \Wone^T\|_\infty = 3.1\times 10^{-4}$
(machine precision effects from eigendecomposition).

\begin{proposition}[Non-Circularity Status of $\Wone$ --- PROVED]\label{prop:W1-circ}
	$\Wone = \CT\cdot\Tt^+$ does not assume RH or the Hilbert--P\'olya postulate.
	However, its normalization constant $\CT$ is the OLS slope in
	$\mu_j\sim\CT/\gamma_{k(j)}$, which is calibrated against the zero ordinates.
	$\Wone$ should therefore not be presented as fully input-free.
	$\Tt$ is explicit and deterministic; the $\gamma_k$ are external
	numerical inputs (Definition~2.2), not conclusions of the model.
	\emph{Proof.}
	$\Tt$ is constructed from prime resonance vectors
	(Definition~2.7), using $\gamma_k$ as numerical inputs $\checkmark$.
	$\Tt^+$ is the Moore--Penrose pseudoinverse --- no postulate $\checkmark$.
	The scalar $\CT$ is computed via OLS regression of $\mu_j$ against
	$1/\gamma_{k(j)}$: this uses $\gamma_k$ as calibration values, but
	does not assume RH or HP. $\square$
\end{proposition}

\textbf{Numerical Result NR~3.4} ($\Wone$ HP-correlation --- failure).
\emph{[Status: NUMERICAL]}

\begin{center}
\begin{tabular}{ccc}
\toprule
$\kappa$ & $\pi(\kappa)$ & $r_2 = \mathrm{corr}(\omega_j,\gamma_{k(j)})$ \\
\midrule
23  & 9  & 0.502 \\
53  & 16 & 0.436 \\
101 & 26 & 0.321 \\
199 & 46 & 0.226 \\
503 & 96 & 0.160 \\
\bottomrule
\end{tabular}
\end{center}

The decreasing trend confirms: $\Wone$ is \textbf{not} an asymptotic
HP-operator.
The structural cause is Gaussian damping: $e^{-\varepsilon^2\gamma_k^2/2}$
in $\Phi$ causes superexponential decay of $\mu_j$
(observed exponent $\alpha\approx -4.88$),
far exceeding the $1/\gamma$ growth pattern required for HP.
Pseudoinversion catastrophically amplifies tail modes.

\textbf{Summary for \S\,3.4:}
$\Wone$ is self-adjoint (proved).
Its construction does not assume RH or the HP postulate, but $\CT$
is calibrated via OLS on the zero ordinates --- $\Wone$ is therefore
not fully input-free (Proposition~\ref{prop:W1-circ}).
$\Wone$ is the natural HP-candidate but fails numerically as
$\kappa$ grows.
This rules out this particular pseudoinverse construction as an HP candidate.
This does not preclude other spectral functions $f(\Tt)$ from being valid
HP-candidates; see Open Problem~6.3.

% -------------------------------------------------------
\section{The Monotonicity Mechanism}

The strategy for bounding $\etaorig$ from below rests on the
interplay of three structural features:
the monotone weight sequence $(c_p)$,
the oscillatory mode structure $(\sin(\gamma_k\log p))$,
and Abel summation.

\subsection{Monotone Weight Structure (Lemma M1) --- [Partial]}

\begin{lemma}[Monotone Weight Structure]\label{lem:M1}
\emph{[Status: PROVED for $p\geq 11$, $K=1$; general case OPEN]}

For $K=1$ (leading term of $V_p(\sigma)$) and $p\geq 11$:
\[
c_p \approx \frac{4\log p}{\sqrt{p}}\quad\text{is strictly monotone
decreasing in } p.
\]
\emph{Proof.}
Differentiate with respect to $p$ (treating $p$ as continuous):
$\frac{d}{dp}[\log p/\sqrt{p}] = (1/2-\log p)/(p\sqrt{p}) < 0$
for $p > e^{1/2}\approx 1.65$.
This holds for all primes $p\geq 2$. $\square$
\end{lemma}

\begin{remark}[Status for $K>1$]
When $K>1$ (i.e., $p^2\leq\kappa$), the weight formula $V_p(\sigma,K)$
is more complex.
The monotonicity $V_p(\sigma,K+1)\geq V_p(\sigma,K)$ is numerically
supported but not formally proved.
The general monotonicity of $c_p$ for all $p$ with $K>1$ is therefore
OPEN.
\end{remark}

\subsection{Oscillatory Mode Structure (Lemma M2) --- [Partial]}

\begin{lemma}[Oscillatory Mode Structure]\label{lem:M2}
\emph{[Partial status --- see below]}

\emph{(M2a) [PROVED]:}
The function $p\mapsto\sin(\gamma_k\log p)$ changes sign in the
intervals
\[
I_m := \bigl[e^{m\pi/\gamma_k},\, e^{(m+1)\pi/\gamma_k}\bigr],
\quad m = 1,2,3,\ldots
\]

\emph{Proof.}
Direct from $\sin(x)=0$ at $x=m\pi$: the argument
$\gamma_k\log p$ crosses $m\pi$ at $p = e^{m\pi/\gamma_k}$. $\square$

\emph{(M2b) [Partial]:}
By the Prime Number Theorem, each interval $I_m$ contains at least
one prime $p\leq\kappa$, for $m$ in a range depending on $\gamma_k$
and $\kappa$.
Full quantitative control for all $m$ is OPEN.
\end{lemma}

\subsection{Abel Summation Principle (Lemma M3) --- [PROVED]}

\begin{lemma}[Abel Summation Principle --- PROVED]\label{lem:M3}
Let $P\geq 1$.
Let $(f_1,\ldots,f_P)$ be monotone non-increasing and non-negative.
Let $(g_1,\ldots,g_P)$ be any sequence with partial sums bounded by $M$:
$|G_m| := |\sum_{i=1}^m g_i|\leq M$ for all $m$.
Then:
\[
\Bigl|\sum_{i=1}^P f_i g_i\Bigr| \leq M\cdot f_1.
\]
\emph{Proof.}
Apply the Abel partial summation formula:
$\sum_{i=1}^P f_i g_i = f_P G_P - \sum_{i=1}^{P-1}(f_{i+1}-f_i)G_i$.
Taking absolute values:
\begin{align*}
\Bigl|\sum_{i=1}^P f_i g_i\Bigr|
&\leq |f_P|\cdot|G_P| + \sum_{i=1}^{P-1}|f_{i+1}-f_i|\cdot|G_i| \\
&\leq Mf_P + M\sum_{i=1}^{P-1}(f_i-f_{i+1})
 \quad\text{[$f$ decreasing]} \\
&= Mf_P + M(f_1-f_P) = Mf_1. \quad\square
\end{align*}
\end{lemma}

\begin{remark}[Application context for Paper~4]
With $f_i = (D^{-1/2}\tilde{c})_{p_i}\propto c_{p_i}/\|a_{p_i}\|$
(monotone decreasing --- Lemma~\ref{lem:M1}) and $g_i = (u_j)_{p_i}$
(oscillatory eigenvector entries --- Lemma~\ref{lem:M2}),
Lemma~\ref{lem:M3} bounds the projection
$|\langle u_j, D^{-1/2}\tilde{c}\rangle|
\leq M_{k(j)}(\kappa)\cdot c_{p_1}/\|a_{p_1}\|$,
provided the partial sums of $u_j$-entries are bounded by
$M_{k(j)}(\kappa)$.
\emph{[The full reduction is structural, not a direct consequence of
Lemma~\ref{lem:M3} alone.]}
\end{remark}

\subsection{Prime Exponential Sum Bound (Theorem EXPLICIT) --- [PROVED]}

\begin{theorem}[Prime Exponential Sum Bound --- PROVED --- no RH used]\label{thm:EXPLICIT}
For fixed $\gamma > 0$ and $\kappa\geq 2$:
\[
M_k(\kappa) := \max_{x\leq\kappa}\Bigl|\sum_{p\leq x}\sin(\gamma_k\log p)\Bigr|
= O\!\left(\frac{\pi(\kappa)}{\gamma_k}\right).
\]
Explicitly:
\[
M_k(\kappa) \leq \left(\frac{2}{\gamma_k} + o\!\left(\frac{1}{\gamma_k}\right)\right)\pi(\kappa),
\quad\kappa\geq 2.
\]
\emph{Proof.}
We establish three auxiliary lemmas.

\textbf{Lemma~T1} (Abel--Stieltjes Representation --- PROVED).
For $\gamma > 0$ and $\kappa\geq 2$:
\[
\sum_{p\leq\kappa}p^{i\gamma}
= A_0(\gamma) + J_\gamma(\kappa) + R_\gamma(\kappa),
\]
where $|A_0(\gamma)|\leq 0.1$;
$J_\gamma(\kappa) := \int_2^\kappa x^{i\gamma}/\log x\,dx$ (main term);
$R_\gamma(\kappa) := \kappa^{i\gamma}E(\kappa)
- i\gamma\int_2^\kappa E(x)x^{i\gamma-1}dx$,
with $E(x) = O(xe^{-c\sqrt{\log x}})$ (error term).
No RH is used; PNT provides the $E(x)$ bound. $\square$

\textbf{Lemma~T2} (Main Term Bound --- PROVED).
\[
|J_\gamma(\kappa)| \leq
\left(\frac{2}{\gamma} + o\!\left(\frac{1}{\gamma}\right)\right)\pi(\kappa).
\]
\emph{Proof.}
The dominant term satisfies
$|\kappa^{1+i\gamma}/((1+i\gamma)\log\kappa)|
= \kappa/(\sqrt{1+\gamma^2}\cdot\log\kappa) \leq \kappa/(\gamma\log\kappa)$.
Since $\pi(\kappa)\sim\kappa/\log\kappa$ (PNT):
$|J_\gamma(\kappa)|\leq(2/\gamma)\pi(\kappa)(1+O(1/\log\kappa))$. $\square$

\textbf{Lemma~T3} (Error Term Control --- PROVED).
For fixed $\gamma > 0$ and $\kappa\to\infty$:
$|R_\gamma(\kappa)| = o(\pi(\kappa))$.

\emph{Proof.}
PNT (Hadamard/de la Vall\'ee-Poussin):
$|E(x)| = O(xe^{-c\sqrt{\log x}})$.
Substitution $u = \sqrt{\log x}$ yields
$\int_2^\kappa e^{-c\sqrt{\log x}}\,dx
= O(\sqrt{\log\kappa}\cdot\kappa\cdot e^{-c\sqrt{\log\kappa}})
= o(\pi(\kappa))$,
since $\sqrt{\log\kappa}\cdot e^{-c\sqrt{\log\kappa}} = o(1/\log\kappa)$.
Combined: $|R_\gamma(\kappa)| = o(\pi(\kappa))$ for fixed $\gamma$. $\square$

\emph{Proof of Theorem~\ref{thm:EXPLICIT}:}
The imaginary part of Lemma~T1 gives
$A_k(\kappa) = \mathrm{Im}[A_0] + \mathrm{Im}[J_{\gamma_k}]
+ \mathrm{Im}[R_{\gamma_k}]$.
Applying T2 and T3:
\[
|A_k(\kappa)| \leq 0.1 + \frac{2}{\gamma_k}\pi(\kappa)(1+o(1)) + o(\pi(\kappa))
= O\!\left(\frac{\pi(\kappa)}{\gamma_k}\right).
\]
The same bound holds for all partial sums (Remark~4.5), so the maximum
$M_k(\kappa)$ satisfies the same bound. $\square$
\end{theorem}

\begin{corollary}
For $k\to\infty$ (with $\gamma_k\to\infty$):
$M_k(\kappa)/\pi(\kappa)\to 0$ uniformly in $\kappa$.
\end{corollary}

\begin{remark}[Goldston--Gonek confirmation \cite{GoldstonGonek}]
An independent formula of Goldston--Gonek yields the same
$O(\pi(\kappa)/\gamma_k)$ bound via mean-value methods.
This provides an independent check.
It provides no improvement for fixed $k$, confirming that the bottleneck
is the joint equidistribution Assumption~A(b).
\end{remark}

\begin{remark}[Methodological note]
The estimate $M_k(\kappa)=O(\pi(\kappa)/\gamma_k)$ is proved
by standard PNT/Abel-summation technology; the analytic method
has precedents in Tur\'an's program on approximative Dirichlet
polynomials \cite{Turan}.
What is specific to the present paper is not the bare
oscillatory-sum bound, but the way this bound enters the
conditional energy-asymmetry argument of \S\,5.
\end{remark}

\begin{remark}[What Theorem~EXPLICIT does NOT prove]
Theorem~\ref{thm:EXPLICIT} does not prove
$M_k(\kappa) = o(\pi(\kappa))$ for \emph{fixed} $k$ as $\kappa\to\infty$.
This stronger statement is equivalent (via Weyl's equidistribution criterion)
to $\zeta(1+in\gamma_k)\neq 0$ for all $n\geq 1$ ---
the 1-dimensional component Assumption~A(a), which is settled
by Hadamard (1896).
The genuinely open bottleneck is the joint equidistribution
Assumption~A(b) (\S\,6).
\end{remark}

% -------------------------------------------------------
\section{Results on $\etaorig > 0$}

\subsection{Numerical Confirmation}

\textbf{Numerical Result NR~5.1} ($\etaorig > 0$ for $\kappa\leq 1009$).
\emph{[Status: NUMERICAL --- $\varepsilon=0.05$, $N=100$]}

Let $\tilde{c} = D^{1/2}c/\sqrt{\Estr}$ be the normalized canonical
weight vector. The Rayleigh quotient identity
\[
\langle\Tren\,\tilde{c},\,\tilde{c}\rangle = 1 - \etaorig
\]
holds to machine precision ($< 2\times 10^{-16}$),
where $\Tren := D^{-1/2}TD^{-1/2}$.

\begin{center}
\begin{tabular}{ccccc}
\toprule
$\kappa$ & $\pi(\kappa)$ & $\langle\Tren\tilde{c},\tilde{c}\rangle$
  & $< 1$? & $\etaorig$ \\
\midrule
23   & 9   & 0.18849 & $\checkmark$ & 0.81151 \\
53   & 16  & 0.33073 & $\checkmark$ & 0.66927 \\
101  & 26  & 0.22803 & $\checkmark$ & 0.77197 \\
199  & 46  & 0.21740 & $\checkmark$ & 0.78260 \\
503  & 96  & 0.21788 & $\checkmark$ & 0.78212 \\
1009 & 169 & 0.18002 & $\checkmark$ & 0.81998 \\
\bottomrule
\end{tabular}
\end{center}

$\etaorig > 0$ for all tested $\kappa\leq 1009$ $\checkmark$.

\textbf{Collective mechanism.}
The dominant contribution to $\langle\Tren\tilde{c},\tilde{c}\rangle
= \sum_j B_j$ comes from mode $j=4$ ($B_4=0.193$),
not from the dominant eigenmode $j=1$ ($B_1=0.040$).
Approximately $92\%$ cancellation occurs at $j=1$ ($\kappa=53$).
The effect is distributed across $\sim 8$ eigenmodes ---
$\etaorig > 0$ is a \textbf{collective spectral effect}, not a simple
$u_1$-orthogonality principle.

\subsection{$\Delta$-Decomposition and Burst Stability}

\begin{definition}[$\Delta$-Decomposition]
With $S(r) := \sum_{k=1}^N e^{-\varepsilon^2\gamma_k^2}\cos(\gamma_k\log r)$:
\[
\Delta = \Dburstvar + \Dcross + \Dstream,
\]
where
\[
\Dburstvar = \sum_{p < q\leq 7}c_p c_q[S(pq)-S(p/q)],
\quad
\Dcross = \sum_{p\leq 7 < q\leq\kappa}c_p c_q[S(pq)-S(p/q)],
\]
\[
\Dstream = \sum_{7 < p < q\leq\kappa}c_p c_q[S(pq)-S(p/q)].
\]
The factorization
$S(pq)-S(p/q) = -2\sum_k e^{-\varepsilon^2\gamma_k^2}
\sin(\gamma_k\log p)\sin(\gamma_k\log q)$
follows from $\cos(A+B)-\cos(A-B) = -2\sin A\sin B$. [PROVED, algebraic]
\end{definition}

\textbf{Numerical Result NR~5.4} ($\Dburstvar$ stability).
\emph{[NUMERICAL]}

\begin{center}
\begin{tabular}{lcccc}
\toprule
Pair & $S(pq)$ & $S(p/q)$ & $c_p c_q$ & Contribution \\
\midrule
$(2,3)$ & 0.9536 & 0.2653 & $\approx 2.408$ & $+1.657$ \\
$(2,5)$ & 0.4110 & 0.6555 & $\approx 1.833$ & $-0.448$ \\
$(2,7)$ & 0.5794 & 0.5812 & $\approx 1.501$ & $-0.003$ \\
$(3,5)$ & 0.9440 & 0.3686 & $\approx 1.531$ & $+0.881$ \\
$(3,7)$ & 0.7193 & 0.5571 & $\approx 1.254$ & $+0.203$ \\
$(5,7)$ & 0.9668 & 0.1135 & $\approx 0.954$ & $+0.814$ \\
\midrule
\textbf{Total} & & & & $\mathbf{\approx +3.11}$ \\
\bottomrule
\end{tabular}
\end{center}

$\Dburstvar\approx 3.11 > 0$, and it is \textbf{$\kappa$-invariant}
for all tested $\kappa > 7$.
The sufficient condition for $\Delta > 0$ is:
\[
|\Dcross + \Dstream| < \Dburstvar \approx 3.11.
\]
This is numerically satisfied for all $\kappa\leq 1009$;
analytic proof is OPEN.

\subsection{Conditional Theorem (Theorem 5.3)}

\begin{theorem}[Conditional and Numerical Results for $\etaorig > 0$]\label{thm:T5}

\textbf{Assumption~A:}

$(a)$~\textbf{[SETTLED --- Hadamard--de la Vall\'ee Poussin 1896]:}
$\zeta(1+in\gamma_k)\neq 0$ for all integers $n\geq 1$
and all $k\leq N$.
This follows directly from Hadamard's theorem, which establishes
$\zeta(1+it)\neq 0$ for all real $t\neq 0$; the values $t = n\gamma_k$
are real and positive, hence covered unconditionally.

$(b)$~The sequences
$\{(\gamma_{k_1}\log p,\,\gamma_{k_2}\log p)\bmod 2\pi\}_{p\leq\kappa,
\kappa\to\infty}$
are jointly equidistributed in $[0,2\pi)^2$
for all $k_1\neq k_2$ with $k_1,k_2\leq N$.

\emph{Structural motivation:}
Part~$(a)$ is precisely what Weyl's equidistribution criterion requires
to make the diagonal sums $\frac{1}{\pi(\kappa)}\sum_p e^{in\gamma_k\log p}$
vanish --- equivalently, it is the non-vanishing condition for $\zeta$ on
$\mathrm{Re}(s)=1$ at the specific values $s=1+in\gamma_k$.
Part~$(b)$ is the natural $2$-dimensional extension required to make the
off-diagonal (cross) terms
$\frac{1}{\pi(\kappa)}\sum_p\sin(\gamma_{k_1}\log p)\sin(\gamma_{k_2}\log p)$
vanish for $k_1\neq k_2$; without it, the cross-term cancellation
in step~(ii) cannot be established.
These two conditions are the minimal structural requirements for the proof
of~(i)--(ii).

\emph{Note:}
Part~$(a)$ is settled by Hadamard (1896):
$\zeta(1+it)\neq 0$ for all real $t\neq 0$ covers $t = n\gamma_k$ directly.
Only Part~$(b)$ is genuinely open.

\begin{remark}[Relation between Part~$(a)$ and Weyl's criterion]
Part~$(a)$ of Assumption~A is settled:
$\zeta(1+it)\neq 0$ for all real $t\neq 0$
(Hadamard--de la Vall\'ee Poussin, 1896), covering $t = n\gamma_k$.
Part~$(b)$ --- the equidistribution of $\{\gamma_k\log p\bmod 2\pi\}$
via Weyl's criterion --- requires the vanishing of
\[
\frac{1}{\pi(\kappa)}\sum_{p\leq\kappa} p^{in\gamma_k} \to 0
\quad\text{as }\kappa\to\infty,\quad\text{for all }n\geq 1,\,k\leq N.
\]
This depends on the \emph{quantitative} behaviour of $\zeta$ near
$s = 1 + in\gamma_k$: while non-vanishing (Part~$(a)$) is known,
the quantitative rate sufficient for Weyl's criterion remains open.
The genuine open content of Assumption~A is therefore Part~$(b)$
alone.
\end{remark}

Under this Assumption:

\begin{enumerate}[label=(\roman*)]
\item \emph{[CONDITIONAL on Assumption~A(b)]}
  The sequence $\{\gamma_k\log p\bmod 2\pi\}_{p\leq\kappa,\kappa\to\infty}$
  is equidistributed modulo $2\pi$ for each fixed $k\leq N$.
\item \emph{[CONDITIONAL on Assumption~A(b) --- Unweighted]}
  The normalized unweighted cross averages vanish:
  \[
  \frac{1}{\pi(\kappa)}\sum_{p\leq\kappa}
  \sin(\gamma_{k_1}\log p)\,\sin(\gamma_{k_2}\log p)\to 0
  \quad\text{for all }k_1\neq k_2.
  \]
\item \emph{[Numerical lower bound --- no assumption needed]}
  $\Delta(\kappa) \geq \Dburstvar \approx 3.11 > 0$ for all tested $\kappa > 7$.
  [NUMERICAL --- $\kappa$-invariant by construction.]
  Numerically, both $\Dcross$ and $\Dstream$ appear to converge individually,
  and their sum tends to $\approx 1.8$; hence
  $\Delta(\kappa) = \Dburstvar + \Dcross + \Dstream
  \to \Delta_\infty \approx 3.11 + 1.8 \approx 4.9 > \Dburstvar$.
  The equidistribution of parts (i)--(ii) implies
  $(\Dcross+\Dstream)/\pi(\kappa)\to 0$, but the absolute values
  converge to a finite positive limit, not to zero.
\item $\etaorig(\kappa)\to\eta_\infty = \Delta_\infty/\Estr(\infty) > 0$.
  [NUMERICAL: $\eta_\infty\approx 0.81$, consistent with
  $\Delta_\infty\approx 4.9$ and $\Estr(\infty)\approx 6.1$.]
  The lower bound from (iii) gives
  $\eta_\infty\geq\Dburstvar/\Estr(\infty)\approx 0.51 > 0$
  unconditionally [NUMERICAL].
\end{enumerate}

\emph{Proof sketch.}
(i) follows from Assumption $(a)$ by Weyl's criterion
(cf.\ \cite{IwaniecKowalski}, \S\,21.1).
(ii) follows from Assumption $(b)$ via
$\sin A\sin B = \tfrac{1}{2}[\cos(A-B)-\cos(A+B)]$.
(iii) is a direct numerical observation: $\Dburstvar$ is $\kappa$-invariant
by construction, and $\Delta(\kappa)\geq\Dburstvar > 0$ for all tested
$\kappa > 7$.
(iv) follows from $\Delta_\infty > 0$ (numerically) and
$\Estr(\infty)\approx 6.1 < \infty$ (proved).
No RH is used. $\square$
\end{theorem}

\begin{remark}[Weighted transfer --- correction note]\label{rem:weighted-transfer}
The equidistribution assumption Assumption~A implies
$\Dstream(\kappa)/\pi(\kappa)\to 0$ and $\Dcross(\kappa)/\pi(\kappa)\to 0$
(Cross/Stream terms are $o(\pi(\kappa))$).
However, numerically $\Dcross + \Dstream$ converges to a finite positive
limit $\approx 1.8$, not to zero.
The earlier formulation ``$\Delta(\kappa)\to\Dburstvar$'' was therefore incorrect:
$\Dburstvar \approx 3.11$ is a \textbf{$\kappa$-invariant lower bound}
for $\Delta(\kappa)$, not the limit.
The correct limit is $\Delta_\infty \approx 4.9 > \Dburstvar$.
The conclusion $\eta_\infty > 0$ follows from the numerical lower bound
$\Delta(\kappa)\geq\Dburstvar > 0$, without requiring the equidistribution
assumption.
\end{remark}

\begin{remark}[Status of the Assumption --- corrected]\label{rem:assumption-status}
Theorem~\ref{thm:T5} requires Assumption~A.
The assumption concerns the line $\mathrm{Re}(s)=1$, distinct from the Riemann
Hypothesis which concerns $\mathrm{Re}(s)=\tfrac{1}{2}$.
The logical relationship between Assumption~A and RH is not analysed in this paper.

\textbf{What is known.}
$\zeta(1+it)\neq 0$ for all real $t\neq 0$
(Hadamard and de la Vall\'ee Poussin, 1896).
This covers \emph{all} values $t = n\gamma_k$ directly ---
$n\gamma_k$ is a positive real number for $n\geq 1$, $k\leq N$,
and Hadamard's theorem is unconditional on the structure of $t$.
Part~$(a)$ of Assumption~A therefore \textbf{requires no proof}: it is a
consequence of Hadamard's theorem.

\textbf{What remains open.}
The equidistribution of $\{\gamma_k\log p\bmod 2\pi\}$ on the circle $[0,2\pi)$.
Via Weyl's criterion, this is equivalent to:
\[
\frac{1}{\pi(\kappa)}\sum_{p\leq\kappa} e^{in\gamma_k\log p} \to 0
\quad\text{as }\kappa\to\infty,\quad\text{for all }n\geq 1,\,k\leq N.
\]
This is a statement about the \emph{average behaviour of primes},
not about individual values of $\zeta$.
It is not implied by Hadamard's theorem alone.
The genuine open content of Assumption~A is Part~$(b)$: equidistribution.

\textbf{Positioning of Part~$(b)$.}
The use of equidistribution conditions on prime-related sequences
as a route toward zero-free regions has precedents in Tur\'an's
program on approximative Dirichlet polynomials \cite{Turan}.
The specific reduction in this paper ---
expressing the energy asymmetry $\etaorig > 0$ as a consequence
of \emph{joint} equidistribution of the multi-parameter family
$\{(\gamma_{k_1}\log p,\,\gamma_{k_2}\log p)\bmod 2\pi\}$
indexed by pairs of zero ordinates --- appears not to be standard.
Assumption~A(b) is materially stronger than classical one-dimensional
equidistribution statements.
\end{remark}

\subsection{Reduction to $\mathrm{Re}(s) = 1$}

\begin{remark}[The line $\mathrm{Re}(s)=1$ vs.\ $\mathrm{Re}(s)=\tfrac{1}{2}$]
Theorem~\ref{thm:T5} requires Assumption~A.
Part $(a)$ translates via Weyl's criterion into a non-vanishing condition:
$\zeta(1+in\gamma_k)\neq 0$ for all $n\geq 1$, $k\leq N$.
This concerns $\mathrm{Re}(s)=1$ --- \textbf{not} $\mathrm{Re}(s)=\tfrac{1}{2}$.

Within the present program, the implication chain is:
\begin{align*}
\mathrm{RH}
&\;\longleftarrow\;
[\text{transfer } \sigma=\tfrac{1}{2}\to\sigma\neq\tfrac{1}{2},\,\text{open}]\\
&\;\longleftarrow\; \etaorig > 0
\;\longleftarrow\; \text{Theorem~\ref{thm:T5}}
\;\longleftarrow\; \text{Assumption~A (open)}.
\end{align*}
The connection to RH requires a quantitative transfer from
$\sigma=\tfrac{1}{2}$ to $\sigma\neq\tfrac{1}{2}$ on the curvature level
--- a step entirely outside the scope of this paper.
\end{remark}

% -------------------------------------------------------
\section{Open Problems}

\begin{openproblem}[Assumption~A --- Main Open Problem]
The proof of Theorem~\ref{thm:T5} requires:

$(a)$~$\zeta(1+in\gamma_k)\neq 0$ for all $n\geq 1$, $k\leq N$.

\emph{Note: Part~$(a)$ alone --- the non-vanishing
$\zeta(1+in\gamma_k)\neq 0$ --- follows directly from
Hadamard's theorem (1896), which establishes
$\zeta(1+it)\neq 0$ for all real $t\neq 0$.
The values $t = n\gamma_k$ are real and positive, hence covered.
The genuine open content of Assumption~A is Part~$(b)$:
the equidistribution of $\{\gamma_k\log p\bmod 2\pi\}$
via Weyl's criterion.}

$(b)$~Joint equidistribution of
$\{(\gamma_{k_1}\log p,\gamma_{k_2}\log p)\bmod 2\pi\}$
for all $k_1\neq k_2$, $k_1,k_2\leq N$.

Known for $(a)$: $\zeta(1+it)\neq 0$ for all real $t\neq 0$
(Hadamard, 1896); this covers $t = n\gamma_k$ directly.
Condition $(b)$ is independent and open.

Possible approaches:
(A) Direct Hadamard product estimate near $s = 1+in\gamma_k$.
(B) Linear independence of $\{\gamma_k\log p\}$ over $\mathbb{Q}$
  (sufficient for $(b)$).
(C) Numerical verification for small $n$, $k$.
\end{openproblem}

\begin{openproblem}[Analytic proof of $\etaorig > 0$ without Assumption]
Prove $\langle\Tren\tilde{c},\tilde{c}\rangle < 1$ for the canonical
weight system without the equidistribution assumption of
Theorem~\ref{thm:T5}.

Current best: $M_k(\kappa) = O(\pi(\kappa)/\gamma_k)$
(Theorem~\ref{thm:EXPLICIT}) bounds individual projections but is
insufficient for the full sum
$\sum_j\lambda_j|\langle u_j, D^{-1/2}\tilde{c}\rangle|^2 < \Estr$.
\end{openproblem}

\begin{openproblem}[HP-operator --- negative result, circularity-free]
Does any circularity-free spectral function $f(\Tt)$ satisfy
$\mathrm{Spec}(f(\Tt))\to\{\gamma_k\}$ as $\kappa\to\infty$?

The natural candidate $\Wone = \CT\cdot\Tt^+$ fails: $r_2\to 0.16$
(\S\,3.4). The structural obstruction is Gaussian damping.

A possible modification: replace $e^{-\varepsilon^2\gamma_k^2/2}$
by $e^{-\varepsilon\gamma_k}$ (slower decay) in $\Phi$.
\end{openproblem}

\begin{openproblem}[Arithmetic constant $\Ceta\approx 0.39$]
Determine the analytic formula for
$\Ceta = \lim_{\kappa\to\infty}\lambda_{\max,\mathrm{ren}}/\pi(\kappa)$.
The candidate $\Ceta = 1/I_N(\varepsilon)$ is falsified.
The correct representation involves a harmonic mean
(analytically open).
\end{openproblem}

\begin{openproblem}[Asymptotics of $\CT(\kappa)/\pi(\kappa)$]
Does $\CT(\kappa)/\pi(\kappa)$ converge as $\kappa\to\infty$?
Numerically: values $2.49, 2.21, 2.17, 1.94, 1.89$ (decreasing) ---
the limit, if it exists, is analytically undetermined.
\end{openproblem}

\begin{openproblem}[Monotone bound]
Is $\etaorig > 0$ a consequence of the Prime Number Theorem alone,
independent of any RH-related hypothesis?

Evidence: $\Dburstvar\approx 3.11 > 0$ depends only on primes
$p\leq 7$ and is $\kappa$-invariant.
Analytically OPEN.
\end{openproblem}

\begin{remark}[Connection to RH --- Geometric Interpretation, not a theorem]
\emph{[Geometric interpretation --- not established as a formal
statement in this paper.]}

The implication chain:
\[
\mathrm{RH} \leftarrow
[\text{transfer open}]
\leftarrow \etaorig > 0
\leftarrow \text{Theorem~\ref{thm:T5}}
\leftarrow \zeta(1+in\gamma_k)\neq 0.
\]
No direct connection between $\etaorig > 0$ and RH is established
in this paper. The missing step is a quantitative transfer from
$\sigma=\tfrac{1}{2}$ to $\sigma\neq\tfrac{1}{2}$ on the curvature level
--- a step entirely outside the scope of this paper and the subject of
future work.
\end{remark}

% -------------------------------------------------------
\section{Fisher-Information Connection (Geometric Interpretation)}

\emph{All statements in this section are clearly marked as geometric
interpretation or [Geometric heuristic --- not used in proof].
None are used in any argument in \S\S\,2--5.}

\subsection{The Weight Vector as Fisher Information}

The canonical weight $c_p = \sqrt{V_p(\tfrac{1}{2})}$ has the structure
of a Fisher information weight at $\sigma=\tfrac{1}{2}$:

\emph{[Geometric interpretation]:}
$V_p(\sigma) = \tfrac{d^2}{d\sigma^2}\log\|M_\sigma f_{\kappa,p}\|^2$
is the log-convexity of the local norm squared ---
analogous to Fisher information for a location family.
The canonical weight $c_p$ assigns larger weight to primes that are
``more informative'' about $\sigma$ at $\sigma=\tfrac{1}{2}$.

\subsection{Energy Interpretation}

\emph{[Geometric interpretation]:}
\[
\Estr = \sum_p V_p(\tfrac{1}{2})\|a_p\|^2
\approx \text{``Fisher-weighted resonance energy''}.
\]
\[
\Delta = \mathrm{Var}\bigl(\sqrt{V_p(\tfrac{1}{2})}\cdot a_p\bigr)
\approx \text{``spectral variance of Fisher-weighted resonance''}.
\]
\emph{[Geometric heuristic --- not a mathematical claim]:}
$\etaorig$ may be interpreted as a ``coding efficiency of primes
for zero localization''.

\subsection{Cram\'er--Rao Analogy}

\emph{[Geometric heuristic --- not used in proof --- not used in any proof]:}

The Cram\'er--Rao inequality reads
$\mathrm{Var}(\hat\theta)\geq 1/\mathcal{I}(\theta)$.
By structural analogy:
$\Estr\longleftrightarrow\mathcal{I}(\tfrac{1}{2})$;
$\Delta\longleftrightarrow\mathrm{Var}$.

\emph{[Speculative analogy --- not a theorem]:}
A Cram\'er--Rao-type inequality
$\etaorig\cdot H_{\mathrm{local}}\geq C$ would be a quantitative
information-theoretic bound; its existence is not established here.
Numerically, $\etaorig\approx 0.67$ and
$H_{\mathrm{local}}\sim 8(\log\kappa)^2\to\infty$ would give a
vanishing lower bound --- consistent but non-informative.
No formal version of this analogy has been proved.

\emph{This section provides intuition for future work, not a proof strategy.}

% -------------------------------------------------------
\section*{Acknowledgments}

The author thanks the reviewers and correspondents whose
comments improved the presentation of this paper.

% -------------------------------------------------------
\section*{Call for Feedback}

This paper is part of an ongoing series toward the Riemann Hypothesis.
All papers are freely available on Zenodo under CC BY 4.0
(\url{https://zenodo.org/search?q=Ulrich+Tehrani}),
and the accompanying code is hosted on GitHub
(\url{https://github.com/utehrani/analysislab-nt}).

The author welcomes critical feedback, corrections, and collaboration.
In particular: analytical approaches to the Weyl equidistribution
of $\{\gamma_k\log p\bmod 2\pi\}$ (Open Problem in \S\,6),
connections to existing results on exponential sums
over primes and zeros of $\zeta(s)$,
and any errors or gaps in the present arguments.

% -------------------------------------------------------
% Bibliography — normative format for the series (v1.4)
% Series papers: Concept-DOI + year only. NO version numbers.
% External refs: standard bibliographic format.
% -------------------------------------------------------
\begin{thebibliography}{99}

\bibitem{BerryKeating}
M.V.~Berry and J.P.~Keating,
``$H=xp$ and the Riemann Zeros'',
in \emph{Supersymmetry and Trace Formulae} (1999), 355--367.

\bibitem{Bombieri}
E.~Bombieri,
``Problems of the Millennium: The Riemann Hypothesis'',
Clay Mathematics Institute, 2000.

\bibitem{BostConnes}
J.-B.~Bost and A.~Connes,
``Hecke algebras, type III factors and phase transitions
with spontaneous symmetry breaking in number theory'',
\emph{Selecta Mathematica, New Series}\ \textbf{1} (1995), 411--457.

\bibitem{Connes}
A.~Connes,
``Trace formula in noncommutative geometry and the zeros
of the Riemann zeta function'',
\emph{Selecta Math.}\ \textbf{5} (1999), 29--106.

\bibitem{GoldstonGonek}
D.A.~Goldston and S.M.~Gonek,
``A note on $S(t)$ and the zeros of the Riemann zeta-function'',
\emph{Bull.\ London Math.\ Soc.}\ \textbf{39} (2007), 482--486;
preprint \emph{arXiv:math/0511092}, 2005.

\bibitem{IwaniecKowalski}
H.~Iwaniec and E.~Kowalski,
\emph{Analytic Number Theory},
AMS Colloquium Publications~53, 2004.

\bibitem{SierraTownsend}
G.~Sierra and P.K.~Townsend,
``Landau levels and Riemann zeros'',
\emph{Physical Review Letters}\ \textbf{101} (2008), 110201.

\bibitem{Turan}
P.~Tur\'an,
``On some approximative Dirichlet-polynomials in the theory of
the zeta-function of Riemann'',
\emph{Danske Videnskabernes Selskab, Mat.-fys.\ Medd.}\ 
\textbf{24}(17) (1948), 36~pp.

% Series papers — Concept-DOI only, no version numbers (v1.4)
% Scoping (v1.7): only papers actually \cite{}'d; no self, no forward
\bibitem{Paper1}
U.~Tehrani,
``A Curvature Decomposition of the Explicit Formula'',
Zenodo, \url{https://doi.org/10.5281/zenodo.19025598}, 2026.

\bibitem{Paper2}
U.~Tehrani,
``From Local Curvature to the Weil Functional'',
Zenodo, \url{https://doi.org/10.5281/zenodo.19106992}, 2026.

\bibitem{Paper3}
U.~Tehrani,
``A Finite-Cutoff Hilbert-Space Model
for Prime--Zero Energy Structure'',
Zenodo, \url{https://doi.org/10.5281/zenodo.19307989}, 2026.

% GitHub code repository — no version tag
\bibitem{Code}
U.~Tehrani,
\emph{Verification scripts for the series},
GitHub, \url{https://github.com/utehrani/analysislab-nt}.

\end{thebibliography}

% -------------------------------------------------------
\appendix

\section{Normative Parameter Table}

\begin{center}
\renewcommand{\arraystretch}{1.3}
\begin{tabular}{lll}
\toprule
Parameter & Value & Definition \\
\midrule
$\kappa$ & 53 (normative benchmark) & Prime cutoff \\
$\varepsilon$ & 0.05 & Gaussian damping \\
$N$ & 100 & Number of zero ordinates \\
$\sigma$ & 0.5 & Reference point \\
$\etaorig(\kappa=53)$ & 0.66927 & Energy asymmetry \\
$\Estr(\infty)$ & $\approx 6.1$ & Limiting structural energy \\
$\eta_\infty$ & $\approx 0.81\pm 0.03$ & Limiting energy asymmetry \\
$\Ceta$ & $\approx 0.39$ & $\lambda_{\max,\mathrm{ren}}/\pi(\kappa)$ \\
$\CT$ (normative) & $\approx 35.3$ ($\kappa=53$)
  & OLS slope $\mu_j\sim\CT/\gamma_{k(j)}$ \\
$\Dburstvar$ & $\approx 3.11$ & $\kappa$-invariant for $\kappa > 7$ \\
$r_1$ & 0.950 ($\kappa=53$)
  & $\mathrm{corr}(\mu_j, 1/\gamma_{k(j)})$ \\
\bottomrule
\end{tabular}
\end{center}

\section{Status Summary}

\begin{center}
\renewcommand{\arraystretch}{1.2}
\begin{tabular}{lll}
\toprule
Result & Section & Status \\
\midrule
$\Tt$ self-adjoint, PSD & \S\,2.2 & PROVED \\
Spectral identity $\lambda_j=\mu_j$ & \S\,2.3 & PROVED \\
$\Tt\neq$ HP-operator & \S\,2.4 & NUMERICAL NEGATIVE RESULT \\
$\Wone$ self-adjoint & \S\,3.4 & PROVED \\
$\Wone$: no RH/HP; $\CT$ OLS-calibrated & \S\,3.4 & PROVED (Prop.~3.8) \\
Lemma~M1 (monotone weights, $K=1$) & \S\,4.1 & PROVED ($p\geq 11$); OPEN ($K>1$) \\
Lemma~M2a (sign changes) & \S\,4.2 & PROVED \\
Lemma~M3 (Abel summation) & \S\,4.3 & PROVED \\
Theorem~T1 (Abel--Stieltjes) & \S\,4.4 & PROVED \\
Lemma~T2 (main term bound) & \S\,4.4 & PROVED \\
Lemma~T3 (error term) & \S\,4.4 & PROVED \\
Theorem~EXPLICIT ($M_k=O(\pi/\gamma_k)$) & \S\,4.4 & PROVED \\
$\mu_j\sim\CT/\gamma_{k(j)}$, $r_1=0.950$ & \S\,3.1 & NUMERICAL \\
Localization $0.45$--$0.99$ & \S\,3.2 & NUMERICAL \\
$\Wone$ fails as HP-operator ($r_2\to 0.16$) & \S\,3.4 & NUMERICAL \\
$\etaorig > 0$ for $\kappa\leq 1009$ & \S\,5.1 & NUMERICAL \\
$\Dburstvar\approx 3.11 > 0$ & \S\,5.2 & NUMERICAL \\
$\Ceta\approx 0.39$ structure & \S\,3.3 & NUMERICAL ($1/I_N$ FALSIFIED) \\
Thm~5.3(i)(ii): equidist.\ vanishing & \S\,5.3 & CONDITIONAL (Assumption~A(b)) \\
Thm~5.3(iii): $\Delta(\kappa)\geq\Dburstvar\approx 3.11 > 0$
  & \S\,5.3 & NUMERICAL \\
Thm~5.3(iv): $\etaorig(\kappa)\to\eta_\infty\approx 0.81 > 0$
  & \S\,5.3 & NUMERICAL \\
Assumption~A(a): $\zeta(1+in\gamma_k)\neq 0$ & \S\,6 & SETTLED (Hadamard 1896) \\
Assumption~A(b): equidistribution & \S\,6 & OPEN \\
$\etaorig > 0$ analytically & \S\,6 & OPEN \\
Closed HP-question for $\Wone$ & \S\,3.4 & CLOSED (falsified) \\
\bottomrule
\end{tabular}
\end{center}

% --------------------------------------------------------
% Footer — normative format (v1.6)
% --------------------------------------------------------
\enlargethispage{3\baselineskip}
\vspace*{\fill}
\begin{minipage}{\linewidth}
\rule{\linewidth}{0.4pt}\smallskip\\
\small\emph{Paper~4 v5 \textperiodcentered\ April~2026
\textperiodcentered\ Pauca sed matura}
\end{minipage}

\end{document}
